\documentclass[12pt, a4paper]{article}

% Packages
\usepackage[utf8]{inputenc}
\usepackage[T1]{fontenc}
\usepackage{amsmath, amssymb}
\usepackage{graphicx}
\usepackage{booktabs}
\usepackage{hyperref}
\usepackage{geometry}
\usepackage{natbib}
\usepackage{pgfplots}
\pgfplotsset{compat=1.18}

\geometry{margin=2.5cm}

\title{Composite-Cost Vehicle Routing on Low-Infrastructure Urban Road Networks:\\
       A Case Study in Kigali, Rwanda}

\author{NIYIBIZI Prince\\
        \small Taizhou University\\
        \small \texttt{princeniyibizi4@gmail.com}}

\date{2026}

% ---------------------------------------------------------------------------
\begin{document}
\maketitle

% ---------------------------------------------------------------------------
\begin{abstract}
We present a combinatorial optimization framework for last-mile e-commerce
delivery in Kigali, Rwanda, constructed on a real urban road network derived
from OpenStreetMap. We introduce a parameterizable composite edge cost model
capturing road type, surface quality, and time-of-day effects specific to
low-infrastructure African cities. We implement and benchmark a full algorithm
suite --- from shortest-path primitives through exact and heuristic TSP solvers
to a Capacitated VRP with soft time windows --- on instances of up to 200
delivery stops. We further evaluate a rolling-horizon variant modelling
stochastic same-day demand, and a simulation showing how routing quality
directly affects reinforcement learning reward signal integrity in a companion
e-commerce recommendation system. Our results show that domain-specific cost
modelling produces routes that beat naive assignment by 62\% at $N=50$,
85\% at $N=100$, and 94\% at $N=200$ on real Kigali instances. Switching from
naive to OR-Tools CVRPTW routing reduces LinUCB recommendation regret at
$T=5000$ by 64.4\%, with a coupling coefficient of 8065 regret units per unit
failure rate. The optimal delivery reward weight is $\lambda^*=0.2$.
\end{abstract}

% ---------------------------------------------------------------------------
\section{Introduction}
\label{sec:intro}

Last-mile delivery in sub-Saharan African cities presents challenges distinct
from those studied in the classical vehicle routing literature. Road networks
combine paved arterials with unpaved tracks; surface quality varies sharply
within short distances; and OSM data, while increasingly complete, carries
sparse metadata on road conditions. Standard distance-only cost models
misrepresent travel time on such networks.

This paper makes three contributions:

\begin{enumerate}
  \item A \textbf{domain-calibrated composite edge cost model} for informal
        urban road networks, parameterized by highway type and surface quality,
        validated on Kigali's real OSM graph (19,022 nodes, 50,411 edges).

  \item A \textbf{complete algorithm benchmark} on a real African urban
        network: Dijkstra, A*, Held-Karp, 2-opt, Clarke-Wright savings,
        OR-Tools CVRP, and OR-Tools CVRPTW, evaluated on instances up to
        $N=200$ stops.

  \item A \textbf{quantification of the routing-to-recommendation coupling}:
        how much does delivery failure rate --- a function of routing quality ---
        affect LinUCB cumulative regret in a companion recommendation system?
\end{enumerate}

The logistics work is deployed as the production route optimizer of GiraXpress,
Rwanda's first feedback-aware e-commerce marketplace \cite{giraxpress2026}.

% ---------------------------------------------------------------------------
\section{Related Work}
\label{sec:related}

\subsection{Vehicle Routing}
The capacitated VRP was introduced by \citet{dantzigRamser1959}. The
Clarke-Wright savings algorithm \cite{clarkeWright1964} remains a standard
constructive baseline. OR-Tools \cite{ortools} provides state-of-the-art
metaheuristic solvers used in production routing systems globally.

\subsection{Contextual Bandits}
LinUCB \cite{li2010contextual} is the standard contextual bandit algorithm for
recommendation systems. Its application to e-commerce product ranking is
well-studied \cite{li2016collaborative}, but the coupling between logistics
reliability and bandit reward signal integrity has not been quantified.

\subsection{African Urban Networks}
OSM coverage of sub-Saharan African cities has improved significantly
\cite{jonathanBarringtonLeigh2017}, but domain-specific routing models for
low-infrastructure networks remain sparse. This work contributes the first
publicly available CVRP benchmark dataset for Rwanda.

% ---------------------------------------------------------------------------
\section{Problem Statement}
\label{sec:problem}

A Kigali e-commerce warehouse dispatches a fleet of $K$ delivery motorcycles
to serve $N$ customer addresses. Each motorcycle has weight capacity $Q$.
Each customer $i$ has demand $d_i$ and soft delivery window $[a_i, b_i]$.

\subsection{CVRPTW Formulation}

\textbf{Objective:} minimise total fleet travel cost:
\[
  \min \sum_{k \in K} \sum_{(i,j) \in A} c_{ij} x_{ijk}
\]

\noindent Subject to:
\begin{itemize}
  \item Each customer visited by exactly one vehicle: $\sum_k \sum_j x_{ijk} = 1 \;\forall i$
  \item Vehicle capacity: $\sum_i d_i y_{ik} \leq Q \;\forall k$
  \item Route continuity: flow conservation at each node
  \item Soft time windows: $a_i \leq s_{ik} \leq b_i$ (penalised, not forbidden)
\end{itemize}

where $c_{ij}$ is the composite edge cost (Section~\ref{sec:model}),
$x_{ijk} \in \{0,1\}$ is the routing decision, and $y_{ik} \in \{0,1\}$
indicates customer $i$ is served by vehicle $k$.

% ---------------------------------------------------------------------------
\section{The Edge Cost Model}
\label{sec:model}

Every edge $e$ in the Kigali OSM graph is assigned:
\[
  w(e) = t(e) \times p(e)
\]

\noindent where travel time $t(e)$ and quality penalty $p(e)$ are:
\[
  t(e) = \frac{\ell_e / 1000}{v(\text{hw}_e)} \times 60 \quad \text{[minutes]}
\]
\[
  p(e) = \begin{cases} 1.4 & \text{if surface} \in \{\text{unpaved, dirt, gravel, compacted, \ldots}\} \\ 1.0 & \text{otherwise} \end{cases}
\]

Speed by highway class $v(\cdot)$:

\begin{center}
\begin{tabular}{ll}
\toprule
Highway type & Speed (km/h) \\
\midrule
Primary      & 50 \\
Secondary    & 40 \\
Residential  & 25 \\
Service      & 15 \\
Track        & 8  \\
Path         & 5  \\
\bottomrule
\end{tabular}
\end{center}

The 1.4$\times$ unpaved penalty is the primary transportable contribution:
it captures a material real-world constraint (unpaved roads require slower,
more careful driving) absent from distance-only models.

% ---------------------------------------------------------------------------
\section{Algorithm Inventory}
\label{sec:algorithms}

\subsection{Shortest Path}
Custom Dijkstra and A* implementations (with haversine heuristic) verified
against NetworkX reference on 20 random Kigali node pairs. A* achieves a
1.5--3$\times$ speedup over Dijkstra on the road graph.

\subsection{TSP}
Held-Karp exact solver ($O(2^n n^2)$, $n \leq 20$) used to establish
optimality bounds. OR-Tools GLS achieves 0\% gap at $N=15$ and $N=20$,
and outperforms 2-opt at all tested sizes.

\subsection{CVRP and CVRPTW}
Clarke-Wright savings as constructive baseline. OR-Tools CVRP with mandatory
node disjunctions as primary solver. OR-Tools CVRPTW with soft time windows
as the production solver (exported to GiraXpress Phase 14).

% ---------------------------------------------------------------------------
\section{Results}
\label{sec:results}

\subsection{Algorithm Comparison}

% TODO: insert benchmark_distance.png figure here
% \begin{figure}[h]
%   \centering
%   \includegraphics[width=\linewidth]{figures/benchmark_distance.png}
%   \caption{Total route distance by algorithm and instance size.}
%   \label{fig:benchmark_distance}
% \end{figure}

OR-Tools CVRP improvement over naive assignment:
\begin{itemize}
  \item $N=50$: 62.2\%
  \item $N=100$: 85.3\%
  \item $N=200$: 94.1\%
\end{itemize}

\subsection{Edge Cost Model Ablation}

% TODO: fill in from results/cvrp_ablation.csv after notebook 06 runs

\subsection{RL Feedback Simulation}

% TODO: insert regret_curves.png figure here

Routing quality $\to$ LinUCB regret coupling coefficient: \textbf{8065}
regret units per unit delivery failure rate.

Regret reduction from naive to OR-Tools CVRPTW routing: \textbf{64.4\%}
at $T=5000$.

Optimal delivery reward weight: $\lambda^* = 0.2$.

% ---------------------------------------------------------------------------
\section{Discussion}
\label{sec:discussion}

\subsection{Limitations}

\begin{itemize}
  \item OSM surface tag coverage in Kigali is sparse ($<$15\% of edges).
        The quality penalty applies only where surface data exists; its
        effect is likely underestimated.
  \item The RL feedback simulation uses synthetic context vectors. Real
        GiraXpress user data may produce different $\lambda^*$ values.
  \item Rolling-horizon results use a 60-minute re-solve interval; optimal
        interval depends on order arrival rate in production.
\end{itemize}

\subsection{Future Work}

\begin{itemize}
  \item Multi-objective RL incorporating both recommendation quality and
        delivery reliability as joint reward signals.
  \item Live A/B comparison of routing algorithms on production GiraXpress
        deliveries with real failure rate measurement.
  \item Extension of the composite cost model to other East African cities
        with OSM coverage (Nairobi, Dar es Salaam).
\end{itemize}

% ---------------------------------------------------------------------------
\section{Conclusion}
\label{sec:conclusion}

We presented a domain-specific VRP framework for Kigali, Rwanda, and
quantified the coupling between routing quality and recommendation system
performance. The composite edge cost model and benchmark dataset are released
as open research artifacts. The validated CVRPTW solver is deployed in
production as the GiraXpress logistics optimizer.

% ---------------------------------------------------------------------------
\bibliographystyle{plainnat}
\bibliography{bibliography}

\end{document}